\documentclass[UTF8]{ctexart}

\makeatletter
\def\input@path{{../}{../../}}
\makeatother
\usepackage{researchnotes}

\title{中止比赛的赌本公平分配}
\author{}
\date{}

\begin{document}
\maketitle

甲、乙两人各拿出150元，共300元作为赌本。双方约定，谁先累计赢下三局，谁就拿走全部赌本。每局没有平局，两人获胜的机会相同，而且前面各局的输赢不会影响后面各局的结果。现在甲已经赢了两局，乙赢了一局，比赛却因外部原因不得不中止。这300元应该怎样公平分配？

\section{问题定义}
\label{sec:points-problem}

\textbf{问题：中止比赛的赌本公平分配。}两名参与者 \(A\) 和 \(B\) 各投入 \(M/2\) 元，形成总额为 \(M>0\) 元的赌本。双方约定进行一系列比赛，每局恰有一名胜者，不存在平局；各局结果相互独立，且每局 \(A\) 获胜的概率均为已知常数 \(p\in(0,1)\)，\(B\) 获胜的概率为 \(1-p\)。比赛采用“先取得 \(s\) 局胜利者获胜”的规则，其中 \(s\) 为正整数：任何一方累计取得第 \(s\) 局胜利时，比赛立即结束，胜者获得全部赌本 \(M\)，败者获得 \(0\)。因此，正常结束的比赛最多进行 \(2s-1\) 局，并不要求打满这些局数。

比赛在完成 \(n=a+b\) 局后，因外部原因永久中止。此时 \(A\) 已赢 \(a\) 局，\(B\) 已赢 \(b\) 局，其中
\[
a,b\in\{0,1,\ldots,s-1\}.
\]
这些条件保证比赛尚未产生最终胜者。假定中止原因不提供关于未来比赛结果的额外信息；若比赛能够继续，后续各局仍遵循上述独立性和胜率假设。双方需要立即分配现有赌本，分配金额分别记为 \(x_A,x_B\)，并要求
\[
x_A\ge 0,\qquad x_B\ge 0,\qquad x_A+x_B=M.
\]

\textcolor{blue}{这里将“公平”明确规定为保持双方在当前信息下的期望所得不变。}设 \(W_A,W_B\) 分别表示假想比赛按原规则继续直至结束时，双方最终获得的金额；它们只可能取 \(0\) 或 \(M\)，且 \(W_A+W_B=M\)。要求提前结算的金额满足
\[
x_A=\mathbb E[W_A\mid\text{当前比分为 }(a,b)],
\qquad
x_B=\mathbb E[W_B\mid\text{当前比分为 }(a,b)].
\]
\textbf{求 \(x_A,x_B\) 关于 \(M,s,a,b,p\) 的表达式。}这里分配的是包含原始投入在内的总赌本；双方净收益分别为 \(x_A-M/2\) 和 \(x_B-M/2\)。

\section{把分配金额转化为最终获胜概率}
\label{sec:points-reduction}

\textcolor{blue}{求解时应当关注双方距离最终胜利还差多少局。}记
\begin{equation}
  r=s-a,\qquad t=s-b,
  \label{eq:points-remaining}
\end{equation}
则 $r,t$ 都是正整数，分别表示 $A,B$ 还需要取得的胜局数。由于每局胜率 $p$ 已知且固定，各局相互独立，中止原因又不提供关于未来结果的额外信息，过去各局的具体排列不会改变后续比赛的分布。因此，对于当前的求解，只需保留 $r,t$ 和 $p$，不必记录前面每一局的胜者。以下概率和期望均在当前比分及中止事实已知的条件下计算。

固定 $p$，用 $Q(r,t)$ 表示从这一状态出发，$A$ 比 $B$ 更早取得所需胜局数的概率；记相应事件为 $E_A$。比赛必在有限局数内产生唯一胜者，因而 $B$ 最终获胜的概率为 $1-Q(r,t)$。由于胜者获得全部赌本，双方最终所得可以写为
\[
  W_A=M\mathbf 1_{E_A},\qquad W_B=M(1-\mathbf 1_{E_A}),
\]
其中 $\mathbf 1_{E_A}$ 是事件 $E_A$ 的\keyterm{示性函数}（indicator function）：事件发生时取 $1$，否则取 $0$。由问题所规定的条件期望分配原则，立即得到
\begin{equation}
  x_A=M Q(r,t),\qquad x_B=M\bigl(1-Q(r,t)\bigr).
  \label{eq:points-allocation}
\end{equation}
\textcolor{blue}{所以全部求解工作归结为计算 $Q(r,t)$。}只要这一概率确定，分配金额便唯一确定，并且自动满足非负性和总额守恒。下面通过补齐未来比赛，将这一概率转化为二项分布的区间概率。

\section{补齐比赛与二项分布}
\label{sec:points-binomial}

\subsection{为什么至多还需要这么多局}

令
\begin{equation}
  N=r+t-1=2s-a-b-1.
  \label{eq:points-horizon}
\end{equation}
如果又进行了 $N$ 局后仍没有胜者，那么 $A$ 至多赢了 $r-1$ 局，$B$ 至多赢了 $t-1$ 局，合计至多进行 $r+t-2=N-1$ 局，与已经进行了 $N$ 局矛盾。因此，真实比赛至多再进行 $N$ 局。

现在预先生成 $N$ 个相互独立的未来比赛结果，每个结果中 $A$ 获胜的概率都是 $p$。\textcolor{blue}{真实比赛仍在某一方达到所需胜局数时立即结束，预先生成的后续结果只作为假想补充，不改变已确定的胜者。}对 $i=1,\ldots,N$，令随机变量 $Y_i$ 在第 $i$ 个未来结果为 $A$ 获胜时取 $1$，否则取 $0$，并令
\[
  K=Y_1+\cdots+Y_N.
\]
于是 $K$ 表示补齐后这 $N$ 局中 $A$ 的总胜局数，服从\keyterm{二项分布}（binomial distribution）：
\begin{equation}
  \mathbb P(K=k)=\binom Nk p^k(1-p)^{N-k},
  \qquad k=0,1,\ldots,N.
  \label{eq:points-binomial-mass}
\end{equation}

\subsection{事件等价与一般分配公式}

\textbf{赌本分配的一般公式。}在第~\ref{sec:points-problem} 节的假设下，令 $r=s-a$、$t=s-b$、$N=r+t-1$，则
\begin{equation}
  Q(r,t)=\sum_{k=r}^{N}\binom Nk p^k(1-p)^{N-k}.
  \label{eq:points-binomial-tail}
\end{equation}
按照条件期望原则，双方应分得
\begin{equation}
  \boxed{
  \begin{aligned}
    x_A&=M\sum_{k=r}^{N}\binom Nk p^k(1-p)^{N-k},\\
    x_B&=M\sum_{k=0}^{r-1}\binom Nk p^k(1-p)^{N-k}.
  \end{aligned}}
  \label{eq:points-general-shares}
\end{equation}
特别地，将 $r,N$ 换回原问题中的参数，有
\begin{equation}
  x_A=M\sum_{k=s-a}^{2s-a-b-1}
    \binom{2s-a-b-1}{k}p^k(1-p)^{2s-a-b-1-k},
  \qquad x_B=M-x_A.
  \label{eq:points-original-parameters}
\end{equation}

若 $A$ 在真实比赛中获胜，补齐后的 $N$ 局中仍至少有 $r$ 个胜局；若 $B$ 获胜，则补齐后至少占据 $t$ 个胜局，因此 $A$ 至多赢 $N-t=r-1$ 局。所以，$A$ 最终获胜恰好对应 $K\ge r$，由二项分布公式~\eqref{eq:points-binomial-mass} 和式~\eqref{eq:points-allocation} 即得上述分配公式。

\section{计算示例与适用范围}
\label{sec:points-examples}

\subsection{原问题的计算}

\begin{example}[三胜制比赛在比分 $2:1$ 时中止]
\label{ex:points-original}
总赌本为 $300$ 元，双方每局胜率相等，先赢三局者取得全部赌本。中止时 $A$ 已赢两局，$B$ 已赢一局，求公平分配金额。
\end{example}

\begin{solve}
由 $r=3-2=1$、$t=3-1=2$，得 $N=2$。补齐法给出
\[
  Q(1,2)=\binom21\left(\frac12\right)^2
          +\binom22\left(\frac12\right)^2
        =\frac34.
\]
因此，双方应分得
\[
  x_A=300\cdot\frac34=225\text{ 元},\qquad
  x_B=300\cdot\frac14=75\text{ 元}.
\]
若保持当前比分和比赛规则不变，但把每局 $A$ 的胜率改为一般的 $p$，由于 $A$ 只差一胜，$B$ 必须连续赢下接下来的两局才能最终获胜，其概率为 $(1-p)^2$。因此，双方应分得
\[
  x_A=300\bigl[1-(1-p)^2\bigr],\qquad x_B=300(1-p)^2.
\]
因此，$225$ 元与 $75$ 元这一具体分配确实依赖于每局双方胜率相等的假设。
\end{solve}

\subsection{双方还需要多局才能获胜的例子}

\begin{example}[四胜制下的不等胜率分配]
\label{ex:points-biased}
双方各投入 $500$ 元，约定先赢四局者获胜。每局 $A$ 获胜的概率为 $p=3/5$，各局独立；中止时 $A$ 已赢一局，$B$ 已赢两局。求双方按条件期望原则应分得的金额。
\end{example}

\begin{solve}
此时 $M=1000$、$r=3$、$t=2$、$N=4$。由二项分布公式，
\[
  Q(3,2)=\binom43\left(\frac35\right)^3\frac25
          +\binom44\left(\frac35\right)^4
        =\frac{216+81}{625}=\frac{297}{625}.
\]
所以 $A$ 分得 $1000\times297/625=475.2$ 元，$B$ 分得 $1000\times328/625=524.8$ 元。\textcolor{blue}{虽然 $A$ 每局的胜率更高，但当前所需胜局数也更多；最终分配同时取决于每局胜率和距离比赛结束的状态，不能只使用其中一个因素。}
\end{solve}

本文的“公平”指第~\ref{sec:points-problem} 节的条件期望原则，即按当前信息保持双方在原赛制下的期望所得；投入相等并不保证期望所得相等。上述公式要求各局独立、胜率 $p\in(0,1)$ 已知且固定，并且中止原因不提供关于未来胜负的额外信息。在这些条件下，只要尚未决出胜者，双方都应分得正金额。若胜率未知或随比赛过程变化，则需先建立相应的统计或状态模型，不能把当前胜局比例直接当作真实胜率套用。

\end{document}
